\subsection*{第9章の索引用語}

\begin{description}
  \item[\term{全体像}] 一要求に関係する層、境界、観測点を一枚にまとめた地図です。
  \item[\term{観測点}] 値、状態、時間、失敗を記録できる境界です。
  \item[\term{span}] 一つの処理区間と属性、開始、終了を表すトレースの単位です。
  \item[\term{trace\_id}] 複数のspanを一要求のトレースへ結び付ける識別子です。
  \item[\term{correlation id}] 異なるログやメッセージを同じ処理へ関連づける識別子です。
  \item[\term{OpenTelemetry}] トレース、メトリクス、ログの計装と伝送を標準化するプロジェクトです。
  \item[\term{SLI}] Service Level Indicatorの略で、サービス水準を測る指標です。
  \item[\term{SLO}] Service Level Objectiveの略で、SLIに対して定める目標です。
  \item[\term{エラーバジェット}] 目標期間内でSLOを満たしつつ許容できる失敗量です。
  \item[\term{プロファイリング}] CPU時間やメモリ割り当ての分布を測る分析です。
  \item[\term{ボトルネック}] 全体の処理量や時間を主に制約する箇所です。
  \item[\term{検証}] 仮説や修正結果を観測によって確かめる行為です。
  \item[\term{再現}] 同じ条件と手順で同じ症状を再び観測することです。
  \item[\term{切り分け}] 観測と反証によって原因候補の範囲を狭めることです。
  \item[\term{DNS委任}] 名前空間の一部を別の権威DNSへ任せる関係です。
  \item[\term{経路障害}] パケットを宛先へ届ける経路が利用できない、または不安定な状態です。
  \item[\term{TLS}] Transport Layer Securityの略で、通信の機密性、完全性、認証を支えます。
  \item[\term{AI品質}] 利用目的に対するモデル出力の正確さ、安全性、安定性などの性質です。
  \item[\term{評価データ}] 品質指標を測るため、学習とは分けて用意する入力と期待結果です。
  \item[\term{インシデント}] サービスへ望ましくない影響が生じ、対応と記録が必要な出来事です。
  \item[\term{ポストモーテム}] インシデント後に事実、影響、原因、改善を非難ではなく学習目的で記録する振り返りです。
  \item[\term{正常系}] 想定した前提と入力で期待する処理が完了する流れです。
  \item[\term{異常系}] 失敗、例外、制限超過など、通常と異なる処理へ進む流れです。
\end{description}
